\documentclass[12pt]{article}
\usepackage{amsmath, amssymb, amsthm}
\usepackage{parskip}
\usepackage{geometry}
\usepackage{enumitem}
\usepackage{fancyhdr}

\geometry{letterpaper, margin=1in}
\setlength{\headheight}{15pt}

% Header setup
\pagestyle{fancy}
\fancyhf{}
\lhead{Math 311 Fall 2026}
\rhead{Jeffrey Lu}
\cfoot{\thepage}

\title{Homework Assignment 4 \\ \large Math 311, Fall 2026 \\ Rutgers University}
\author{}
\date{}

\begin{document}

\maketitle

\textbf{Exercise 2.5.1}
\begin{enumerate}[label=(\alph*)]
    \item By Bolzano Wierstrass, every bounded sequence has a convergent subsequence.
          So if a sequence has a bounded subsequence, then that subsequence must also contain a convergent subsequence,
          which would be a subsequence of the main sequence as well.

    \item Define
          \[
              x_{2n-1} = \frac{1}{n+1}, \qquad x_{2n} = 1 - \frac{1}{n+1}.
          \]
          Then \( x_n \neq 0,1 \) but
          \[
              x_{2n-1} \to 0, \qquad x_{2n} \to 1.
          \]

    \item If we combine the sequences
          \[
              (1), (1, \frac{1}{2}), (1,\frac{1}{2},\frac{1}{3}), \dots
          \]
          together to form
          \[
              1, 1, \frac{1}{2}, 1, \frac{1}{2}, \frac{1}{3}, \dots
          \]
          For each $n \in \mathbb{N}$, we can get a constant subsequence of \( \frac{1}{n} \) by
          choosing the nth element of each "block".

    \item This is impossible, since the sequence formed by the infinite set ordered, eventually converges to 0, but 0 is not in the set.


\end{enumerate}

\textbf{Exercise 2.5.2}
\begin{enumerate}[label=(\alph*)]
    \item True, consider the proper subsequence \( (x_2, x_3, \dots) \), all the elements except the first.
          Then this subsequence converges by assumption. Adding \( x_1 \) to the front does not affect the tail, so \( (x_n) \)
          converges as well.
    \item True, if \( (x_n) \) converges, then every subsequence also converges. Thus the contrapositive means that if we a divergent subsequence, then the \( (x_n) \) diverges.
    \item True, by Bolzano Weierstrass, if \( (x_n) \) is bounded,
          then it must have a convergent subsequence, say \( (x_{nk} \to L) \).
          Then since \( (x_n) \) is divergent, there exists \( \epsilon > 0 \)
          such that infinitely many terms satisfy \( \vert x_n - L \vert \geq \epsilon > 0 \).
          Using these terms, we can construct another bounded subsequence that never converges to L.
          Since this subsequence is bounded, we can reapply Bolzano Weierstrass
          and find a convergent sequence that converges to
          another limit that cannot be \( L \) since \( \epsilon > 0 \).

    \item Suppose \( (x_n) \) is increasing WLOG and contains a convergent subsequence \( x_{nk} \to L \).
          Since the subsequence is convergent, it is also bounded by some upper limit say \( B \).
          Then for every \( n \in \mathbb{N} \), choose \( nk \geq n \) so
          \[
              x_n \leq x_{nk} < B.
          \]
          Thus the \( (x_n) \) is also bounded above by \( B \). Thus, by the monotone convergence theorem, \( (x_n) \to L\).
\end{enumerate}

\textbf{Exercise 2.5.5}
Assume \( (a_n) \) is bounded sequence where every convergent subsequence converge ot the same limit \( a \in \mathbb{R} \).
Since \( a_n \) is bounded, exercise 2.5.2c says that if \( a_n \) diverges, then it's subsequences converge to different limit.
This contradicts our assumption that all convergent subsequences converge to \( a \). Thus \( a_n \) is convergent.
Then all its subsequences must all be convergent, and they all converge to \( a \). Therefore \( a_n \to a \).

\textbf{Exercise 2.5.9}
Let \( (a_n) \) be a bounded sequence, and define
\[
    S = \{x \in \mathbb{R} : x < a_n \text{ for infinitely many terms } a_n\}.
\]
Since \( (a_n) \) is bounded,  there exists some \( B \in \mathbb{R} \) such that
\[
    -B \leq a_n \leq B \text{ for every } n \in \mathbb{N}.
\]
Then \( -B-1 \in S\), so \( S \neq \varnothing  \) and every element \( x \in S \) must be less than \( B \) since
\[
    x < a_n \leq B.
\]
That is, \( S \) is bounded above, so by the axiom of completeness, S must have a least upper bound \( s =\sup S \).
We construct a subsequence \( (a_{nk}) \) out of \( S \) by recursively choosing increasing elements.
Then \( a_{nk} \) is bounded as well, so by the monotone convergence theorem, \( a_{nk} \) converges to some limit \( L \).

For each \( k \in \mathbb{N} \), since \( s - \frac{1}{k} < s \), it follows that
\( s - \frac{1}{k} \) is not an upper bound for \( S \).
So there exists \( t \in S \) such that
\[
    s - \frac{1}{k}  < t.
\]

Then by definition of \( S \), there are infinitely many terms that satisfy \( a_n > t \).
So infinitely many terms also satisfy
\[
    a_n > s - \frac{1}{k}.
\]
While, only finitely many terms satisfy
\[
    a_n > s + \frac{1}{k}.
\]
We can then recursively choose indices \( n_1 < n_2 < \dots  \) so that
\[
    s - \frac{1}{k} < a_{nk} \le s + \frac{1}{k}.
\]
That is,
\[
    \vert a_{nk} - s \vert \leq  \frac{1}{k}.
\]
Given \( \epsilon > 0 \), choose \( K \in \mathbb{N} \) such that \( \frac{1}{K} < \epsilon \).
Then, for every \( k \geq K \) we have
\[
    \vert a_{nk} - s \vert \leq \frac{1}{k} \leq \frac{1}{K} < \epsilon
\]
Therefore,
\( a_{nk} \to s = \sup S.
\)

\textbf{Exercise 2.6.2}
\begin{enumerate}[label=(\alph*)]
    \item Let
          \[
              x_n = \frac{(-1)^n}{n}.
          \]
          Then \( x_n \to 0 \), so it is cauchy.
          However, it is not monotone since it alternates between positive and negative values.
    \item Impossible.
          Every cauchy sequence is convergent, every convergent sequence is bounded,
          and every subsequence of a bounded sequence is bounded.
    \item Impossibe. Since the subsequence is cauchy, it must also be convergent, and so it must also be bounded.
          Since the sequence is monotone, if it has a bounded subsequence, it must be bounded as well.
          Then the sequence is bounded and monotone, so by the monotoen convergent theorem, it must be convergent.

    \item Consider the sequence \( (x_n) \) defined by
          \[
              x_{2n} = 0, \qquad  x_{2n-1} = n.
          \]
          Then the sequence is unbounded since \( x_{2n-1} \to \infty \),
          but its even subsequence \( x_{2n} \to 0 \), and so \( x_n \) contains a cauchy.

\end{enumerate}


\textbf{Exercise 2.6.4}
\begin{enumerate}[label=(\alph*)]
    \item Yes. Let \( \epsilon > 0 \).
          Since \( a_n \) and \( b_n \) are cauchy, choose \( N \) such that for all \( m,n \geq N \),
          \[
              \vert a_n - a_m \vert < \frac{\epsilon}{2}, \qquad \vert b_n - b_m \vert < \frac{\epsilon}{2}.
          \]
          Then for every \( m,n \geq N \) we have
          \[
              \begin{split}
                  \vert c_n - c_m \vert & = \vert \vert a_n - b_n \vert - \vert a_m - b_m \vert  \vert \\
                                        & \leq \vert (a_n - b_n) - (a_m - b_m) \vert                   \\
                                        & = \vert (a_n - a_m) + (b_m - b_n) \vert                      \\
                                        & \leq \vert a_n - a_m \vert + \vert b_n - b_m \vert           \\
                                        & < \frac{\epsilon}{2} + \frac{\epsilon}{2}                    \\
                                        & = \epsilon.
              \end{split}
          \]
          Thus \( c_n \) is cauchy.
    \item Not cauchy. Consider \( (a_n) = 1 \).
          Then \( a_n \to 1 \) and so is cachy.
          But \( c_n = (-1)^n \) which diverges and thus is not cauchy.

    \item Not cauchy. Consider \( a_n = \frac{(-1)^n}{n} \).
          Since \( a_n \to 0 \), \( a_n \) is cauchy. But
          \[
              c_n = \begin{cases}
                  0,  & n \text{ even} \\
                  -1, & n \text{ odd}  \\
              \end{cases}
          \]
          Thus, the sequence is alternating and never converges. Therefore \( c_n \) is not cauchy.


\end{enumerate}






\end{document}